\documentclass[11pt,a4paper]{article}
\usepackage[T1]{fontenc}
\usepackage[utf8]{inputenc}
\usepackage{amsmath,amssymb,amsthm}
\usepackage[margin=1in]{geometry}
\usepackage[colorlinks=true,linkcolor=blue,urlcolor=blue]{hyperref}
\theoremstyle{plain}
\newtheorem{theorem}{Theorem}
\newtheorem{lemma}[theorem]{Lemma}
\newtheorem{proposition}[theorem]{Proposition}
\newtheorem{corollary}[theorem]{Corollary}
\theoremstyle{definition}
\newtheorem{remark}[theorem]{Remark}
% A quoted conjecture can be a single hundred-term formula whose terms are thirty-digit
% numbers. TeX cannot hyphenate those, so without a little stretch every such quote runs off
% the page; the linked-a-file papers are all of that shape.
\setlength{\emergencystretch}{3em}

\title{The empirical closed form for OEIS A227261, proved}
\author{Adrian Perez Fontelles\\ \small Independent researcher}
\date{1 September 2026}
\begin{document}
\maketitle

\begin{abstract}
OEIS A227261 carries an empirical closed form for $a(n)$ --- not a recurrence relating
consecutive terms but an explicit formula. It is true, and it is decidable rather than
empirical. The entry counts arrays subject to a condition on a bounded window of consecutive lines, so the lines of an array are the
vertices of a finite digraph, an array is a walk in it, and $a(n)$ is a walk count on
$S=235$ vertices. Any function of the form $\sum_j p_j(n)\lambda_j^{\,n}$ satisfies the
monic linear recurrence whose characteristic polynomial is
$\prod_j(x-\lambda_j)^{\deg p_j+1}$; here that polynomial is $q(x)=(x-1)^{20}$, of degree
$20$. Two things are then checked exactly: that the walk count satisfies the same
recurrence from some index on, which the Cayley--Hamilton theorem reduces to a finite
computation, and that the two sequences agree at $20$ consecutive indices past that
point. A monic recurrence determines every later term from its predecessors, so agreement
there is agreement for good.
\end{abstract}

\noindent\small 2020 Mathematics Subject Classification. 05A15, 11B37, 15A18.\normalsize

\section{The conjecture}

OEIS A227261 is ``Number of n X 4 0,1 arrays indicating 2 X 2 subblocks of some larger (n+1) X 5 binary array having a sum of two or less, with rows and columns of the latter in lexicographically nondecreasing order.''. It has offset $1$ and begins
\[
5, 50, 353, 2201, 11932, 57146, 244818, 951917,\ \dots
\]
The entry states, as a conjecture contributed by its author and never marked settled:
\begin{quote}\raggedright\small
Empirical: a(n) = (1/1689515283456000)*n\^{}19 + (43/1067062284288000)*n\^{}18 + (1/1097800704000)*n\^{}17 + (1/24908083200)*n\^{}16 + (2161/2615348736000)*n\^{}15 + (7723/10461394944000)*n\^{}14 + (76931/201180672000)*n\^{}13 + (38351/48283361280)*n\^{}12 - (9767641/402361344000)*n\^{}11 + (9434683/6967296000)*n\^{}10 - (3668820107/402361344000)*n\^{}9 - (5124736597/80472268800)*n\^{}8 + (150774753091/72648576000)*n\^{}7 - (42587741190289/3923023104000)*n\^{}6 - (18276563116219/163459296000)*n\^{}5 + (6577553955799/3113510400)*n\^{}4 - (743180805841567/51459408000)*n\^{}3 + (63341814516853/1286485200)*n\^{}2 - (326223749821/4476780)*n + 17722 for n>9.
\end{quote}
As of the ``Last modified'' line on the live entry (Jun 26 2025 12:39:40, revision 6) this is still
recorded as empirical, and nothing on the entry records it as proved. Write
\begin{multline*}
f(n)\;=\;\frac{n^{19}}{1689515283456000} + \frac{43 n^{18}}{1067062284288000} \\
+ \frac{n^{17}}{1097800704000} + \frac{n^{16}}{24908083200} \\
+ \frac{2161 n^{15}}{2615348736000} + \frac{7723 n^{14}}{10461394944000} \\
+ \frac{76931 n^{13}}{201180672000} + \frac{38351 n^{12}}{48283361280} \\
- \frac{9767641 n^{11}}{402361344000} + \frac{9434683 n^{10}}{6967296000} \\
- \frac{3668820107 n^{9}}{402361344000} - \frac{5124736597 n^{8}}{80472268800} \\
+ \frac{150774753091 n^{7}}{72648576000} - \frac{42587741190289 n^{6}}{3923023104000} \\
- \frac{18276563116219 n^{5}}{163459296000} + \frac{6577553955799 n^{4}}{3113510400} \\
- \frac{743180805841567 n^{3}}{51459408000} + \frac{63341814516853 n^{2}}{1286485200} \\
- \frac{326223749821 n}{4476780} + 17722.
\end{multline*}

\section{The count is a walk count}

The entry's defining condition is local: it is a condition on a bounded window of consecutive lines. Take as vertices the
admissible configurations of such a window and put an edge from $u$ to $v$ when $v$ may
follow $u$, which the condition decides by inspection. An admissible array is then exactly a
walk, so with $M$ the adjacency matrix of that digraph on $S=235$ vertices, and $u,w$ the
vectors recording which windows may start and end an array,
\[
a(n)\;=\;u^{\!\top}M^{\,g(n)}w
\]
for the linear function $g$ that the entry's shape supplies. In particular $a$ is
$C$-finite: it satisfies some monic integer linear recurrence, though not necessarily a short
one.

\section{A closed form is a recurrence}

\begin{lemma}\label{lem:ann}
Let $f(m)=\sum_{j}p_j(m)\lambda_j^{\,m}$ with the $\lambda_j$ distinct and the $p_j$
polynomials, and put $q(x)=\prod_j(x-\lambda_j)^{\deg p_j+1}$. Then $q$ applied to the shift
operator annihilates $f$: writing $q(x)=x^{r}-\sum_{i=1}^{r}c_ix^{r-i}$,
\[
f(m)\;=\;\sum_{i=1}^{r}c_i\,f(m-i)\qquad\text{for every }m.
\]
\end{lemma}

\begin{proof}
The shift operator $E$ acts on $m\mapsto\lambda^{m}p(m)$ as
$(E-\lambda)\bigl(\lambda^{m}p(m)\bigr)=\lambda^{m+1}\bigl(p(m+1)-p(m)\bigr)$, and
$p(m+1)-p(m)$ has degree one less than $p$. So $(E-\lambda)^{\deg p+1}$ kills
$\lambda^{m}p(m)$, and the factors of $q$ commute.
\end{proof}

Here $q(x)=(x-1)^{20}$, of degree $20$; the closed form is a polynomial in $n$, so this is $(x-1)^{20}$, and the recurrence it encodes is
\begin{equation}\label{eq:rec}
a(n)\;=\;20\,a(n-1) - 190\,a(n-2) + 1140\,a(n-3) - 4845\,a(n-4) + 15504\,a(n-5) - 38760\,a(n-6) + 77520\,a(n-7) - 125970\,a(n-8) + 167960\,a(n-9) - 184756\,a(n-10) + 167960\,a(n-11) - 125970\,a(n-12) + 77520\,a(n-13) - 38760\,a(n-14) + 15504\,a(n-15) - 4845\,a(n-16) + 1140\,a(n-17) - 190\,a(n-18) + 20\,a(n-19) - a(n-20).
\end{equation}

\section{The proof}

\begin{lemma}\label{lem:crit}
With $q$ as above, $a$ satisfies \eqref{eq:rec} for all large $n$ if and only if
$u^{\!\top}M^{\,j}q(M)w=0$ for every $j\ge0$, and it is enough to check $j<S$.
\end{lemma}

\begin{proof}
Substituting the walk expression, the residual $a(n)-\sum_ic_ia(n-i)$ equals
$u^{\!\top}M^{\,g(n)-r}q(M)w$ up to the fixed linear reparametrisation $g$. For the bound,
$M^{S}$ is an integer combination of $I,M,\dots,M^{S-1}$ by Cayley--Hamilton, so the values
$u^{\!\top}M^{j}q(M)w$ satisfy the monic recurrence given by the characteristic polynomial of
$M$; $S$ consecutive zeros therefore force all later ones, and the last nonzero value pins the
threshold exactly.
\end{proof}

\begin{theorem}
$a(n)=f(n)$ for every $n\ge10$.
\end{theorem}

\begin{proof}
The vector $w'=q(M)w$ was formed by $20$ matrix--vector products in exact integer
arithmetic, and $u^{\!\top}M^{j}w'$ evaluated for $j=0,1,\dots$, again exactly. Those integers
vanish from the point corresponding to $n=30$ onwards, and the run was continued until the
working vector was identically zero, which settles every larger $j$ at once. So $a$ satisfies
\eqref{eq:rec} for every $n>29$. By Lemma~\ref{lem:ann} $f$ satisfies \eqref{eq:rec}
identically. Both were then evaluated in exact arithmetic at the $20$ consecutive indices
$n=30,\dots,49$ and agree at each. A monic recurrence of order $20$
determines $a(m)$ from $a(m-1),\dots,a(m-20)$, and for every $m\ge50$ both
$m>29$ and $m-20\ge30$ hold, so induction on $m$ gives $a(m)=f(m)$ for every
$m\ge30$.
Below $n=30$ the recurrence argument gives nothing, since \eqref{eq:rec} is only
established for $a$ from $n=30$ on. The remaining indices $n=10,\dots,29$ were
therefore checked one at a time, directly against the walk count in exact integer arithmetic,
and agree.
\end{proof}

The range is tight in the sense that was checked: at $n=9$ the closed form and the sequence differ, so no larger range is available.

\section{Verification}

Four checks, each independent of the algebra above.

First, the walk model reproduces every one of the $22$ terms the entry publishes,
exactly, in integer arithmetic. This is what ties the model to the entry: the model is built
by reading the entry's English, and a misreading gives a different digraph and different
counts.

Second, the closed form was evaluated in exact rational arithmetic --- not floating point ---
at every published term from $n=10$ on, and agrees with the entry's own DATA at each.

Third, the annihilation test of Lemma~\ref{lem:crit} was carried out exactly, and the model
and the closed form were then compared directly at sixty consecutive indices beyond
$n=10$, far past where the proof already settles the matter.

Fourth, the closed form was deliberately perturbed --- by a constant and by a term linear in
$n$ --- and each perturbed form was rejected by the same comparison, so the test is not one
that anything would pass.

\begin{thebibliography}{9}
\bibitem{oeis} The OEIS Foundation, \emph{The On-Line Encyclopedia of Integer
Sequences}, \url{https://oeis.org}, 2026. Sequence A227261.
\bibitem{stanley} R.~P.~Stanley, \emph{Enumerative Combinatorics}, Volume 1, 2nd ed.,
Cambridge University Press, 2011. (Transfer-matrix method, Section 4.7; $C$-finite sequences,
Section 4.1.)
\bibitem{kp} M.~Kauers and P.~Paule, \emph{The Concrete Tetrahedron}, Springer, 2011.
(Closed forms and linear recurrences with constant coefficients.)
\bibitem{horn} R.~A.~Horn and C.~R.~Johnson, \emph{Matrix Analysis}, 2nd ed., Cambridge
University Press, 2013.
\end{thebibliography}

\end{document}
